% --- 题目 ---
\problem[P224-2 (25-3-仅数学三)]{设总体 $\displaystyle X$ 的分布函数为 $\displaystyle F(x)$, $\displaystyle X_1, X_2, \dots, X_n$ 为来自总体 $\displaystyle X$ 的简单随机样本, 样本的经验分布函数为 $\displaystyle F_n(x)$. 对于给定的 $\displaystyle x$ ($\displaystyle 0 < F(x) < 1$), $\displaystyle D[F_n(x)] = ( \quad )$
\begin{tasks}(2)
  \task $\displaystyle F(x)[1-F(x)]$.
  \task $\displaystyle [F(x)]^2$.
  \task $\displaystyle \frac{1}{n} F(x)[1-F(x)]$.
  \task $\displaystyle \frac{1}{n} [F(x)]^2$.
\end{tasks}}
\ansat{369；【十年真题】 - 考点一：抽样分布 - 2}

% --- 题目 ---
\problem[P224-3 (23-1,3)]{设 $\displaystyle X_1, X_2, \dots, X_n$ 为来自总体 $\displaystyle N(\mu_1, \sigma^2)$ 的简单随机样本, $\displaystyle Y_1, Y_2, \dots, Y_m$ 为来自总体 $\displaystyle N(\mu_2, 2\sigma^2)$ 的简单随机样本, 且两样本相互独立. 记 $\displaystyle \bar{X} = \frac{1}{n} \sum_{i=1}^n X_i, \bar{Y} = \frac{1}{m} \sum_{i=1}^m Y_i$, $\displaystyle S_1^2 = \frac{1}{n-1} \sum_{i=1}^n (X_i - \bar{X})^2, \ S_2^2 = \frac{1}{m-1} \sum_{i=1}^m (Y_i - \bar{Y})^2,$ 则 ( \quad )
\begin{tasks}(2)
  \task $\displaystyle \frac{S_1^2}{S_2^2} \sim F(n, m)$.
  \task $\displaystyle \frac{S_1^2}{S_2^2} \sim F(n-1, m-1)$.
  \task $\displaystyle \frac{2S_1^2}{S_2^2} \sim F(n, m)$.
  \task $\displaystyle \frac{2S_1^2}{S_2^2} \sim F(n-1, m-1)$.
\end{tasks}}
\ansat{369；【十年真题】 - 考点一：抽样分布 - 3}

% --- 题目 ---
\problem[P224-4 (18-3)]{设 $\displaystyle X_1, X_2, \dots, X_n (n \ge 2)$ 为来自总体 $\displaystyle N(\mu, \sigma^2)$ ($\displaystyle \sigma > 0$) 的简单随机样本. 令 $\displaystyle \bar{X} = \frac{1}{n} \sum_{i=1}^n X_i, \ S = \sqrt{\frac{1}{n-1} \sum_{i=1}^n (X_i - \bar{X})^2}, S^* = \sqrt{\frac{1}{n} \sum_{i=1}^n (X_i - \mu)^2},$
则 ( \quad )
\begin{tasks}(2)
  \task $\displaystyle \frac{\sqrt{n}(\bar{X} - \mu)}{S} \sim t(n)$.
  \task $\displaystyle \frac{\sqrt{n}(\bar{X} - \mu)}{S} \sim t(n-1)$.
  \task $\displaystyle \frac{\sqrt{n}(\bar{X} - \mu)}{S^*} \sim t(n)$.
  \task $\displaystyle \frac{\sqrt{n}(\bar{X} - \mu)}{S^*} \sim t(n-1)$.
\end{tasks}}
\ansat{369；【十年真题】 - 考点一：抽样分布 - 4}

% --- 题目 ---
\problem[P224-4 (18-3)]{设 $\displaystyle X_1, X_2, \dots, X_n (n \ge 2)$ 为来自总体 $\displaystyle N(\mu, \sigma^2)$ ($\displaystyle \sigma > 0$) 的简单随机样本. 令 $\displaystyle \bar{X} = \frac{1}{n} \sum_{i=1}^n X_i, \ S = \sqrt{\frac{1}{n-1} \sum_{i=1}^n (X_i - \bar{X})^2}, S^* = \sqrt{\frac{1}{n} \sum_{i=1}^n (X_i - \mu)^2},$
则 ( \quad )
\begin{tasks}(2)
  \task $\displaystyle \frac{\sqrt{n}(\bar{X} - \mu)}{S} \sim t(n)$.
  \task $\displaystyle \frac{\sqrt{n}(\bar{X} - \mu)}{S} \sim t(n-1)$.
  \task $\displaystyle \frac{\sqrt{n}(\bar{X} - \mu)}{S^*} \sim t(n)$.
  \task $\displaystyle \frac{\sqrt{n}(\bar{X} - \mu)}{S^*} \sim t(n-1)$.
\end{tasks}}
\ansat{369；【十年真题】 - 考点一：抽样分布 - 4}

% --- 题目 ---
\problem[P224-1 (23-3-仅数学三)]{设 $\displaystyle X_1, X_2$ 为来自总体 $\displaystyle N(\mu, \sigma^2)$ 的简单随机样本, 其中 $\displaystyle \sigma (\sigma > 0)$ 是未知参数. 记 $\displaystyle \hat{\sigma} = a|X_1 - X_2|$, 若 $\displaystyle E(\hat{\sigma}) = \sigma$, 则 $\displaystyle a = ( \quad )$
\begin{tasks}(4)
  \task $\displaystyle \frac{\sqrt{\pi}}{2}$.
  \task $\displaystyle \frac{\sqrt{2\pi}}{2}$.
  \task $\displaystyle \sqrt{\pi}$.
  \task $\displaystyle \sqrt{2\pi}$.
\end{tasks}}
\ansat{370；【十年真题】 - 考点二：统计量的数字特征 - 1}

% --- 题目 ---
\problem[P224-3 (24-3-仅数学三)]{设总体 $\displaystyle X$ 服从 $\displaystyle [0, \theta]$ 上的均匀分布, 其中 $\displaystyle \theta \in (0, +\infty)$ 为未知参数. $\displaystyle X_1, X_2, \dots, X_n$ 是来自总体 $\displaystyle X$ 的简单随机样本, 记 $\displaystyle X_{(n)} = \max\{X_1, X_2, \dots, X_n\}, T_c = cX_{(n)}$.
\begin{enumerate}
\item[(1)] 求 $\displaystyle c$, 使得 $\displaystyle E(T_c) = \theta$;
\item[(2)] 记 $\displaystyle h(c) = E[(T_c - \theta)^2]$, 求 $\displaystyle c$ 使得 $\displaystyle h(c)$ 最小.
\end{enumerate}}
\ansat{370；【十年真题】 - 考点二：统计量的数字特征 - 3}

% --- 题目 ---
\problem[P227-7 (99-3)]{在天平上重复称量一重为 $\displaystyle a$ 的物品, 假设各次称量结果相互独立且都服从正态分布 $\displaystyle N(a, 0.2^2)$. 若以 $\displaystyle \bar{X}_n$ 表示 $\displaystyle n$ 次称量结果的算术平均值, 则为使 $\displaystyle P\{|\bar{X}_n - a| < 0.1\} \ge 0.95$, $\displaystyle n$ 的最小值应不小于自然数 \underline{\hspace{4em}}.
(注: 标准正态分布函数值 $\displaystyle \Phi(1.96) = 0.975$)}
\ansat{370；【真题精选】 - 考点一：抽样分布 - 7}

% --- 题目 ---
\problem[P227-8 (98-3)]{设 $\displaystyle X_1, X_2, X_3, X_4$ 为来自正态总体 $\displaystyle N(0, 2^2)$ 的简单随机样本;
\[ X = a(X_1 - 2X_2)^2 + b(3X_3 - 4X_4)^2, \]
其中 $\displaystyle a, b \neq 0$. 则当 $\displaystyle a = \underline{\hspace{4em}}, b = \underline{\hspace{4em}}$ 时, 统计量 $\displaystyle X$ 服从 $\displaystyle \chi^2$ 分布, 其自由度为 \underline{\hspace{4em}}.}
\ansat{370；【真题精选】 - 考点一：抽样分布 - 8}

% --- 题目 ---
\problem[P227-2 (11-3)]{设总体 $\displaystyle X$ 服从参数为 $\displaystyle \lambda (\lambda > 0)$ 的泊松分布, $\displaystyle X_1, X_2, \dots, X_n (n \ge 2)$ 为来自总体 $\displaystyle X$ 的简单随机样本, 则对应的统计量 $\displaystyle T_1 = \frac{1}{n} \sum_{i=1}^n X_i, T_2 = \frac{1}{n-1} \sum_{i=1}^{n-1} X_i + \frac{1}{n} X_n$, 有 ( \quad )
\begin{tasks}(2)
  \task $\displaystyle ET_1 > ET_2, DT_1 > DT_2$.
  \task $\displaystyle ET_1 > ET_2, DT_1 < DT_2$.
  \task $\displaystyle ET_1 < ET_2, DT_1 > DT_2$.
  \task $\displaystyle ET_1 < ET_2, DT_1 < DT_2$.
\end{tasks}}
\ansat{370；【真题精选】 - 考点二：统计量的数字特征 - 2}

% --- 题目 ---
\problem[P227-5 (08-1,3)]{设 $\displaystyle X_1, X_2, \dots, X_n$ 是总体 $\displaystyle N(\mu, \sigma^2)$ 的简单随机样本. 记 $\displaystyle \bar{X} = \frac{1}{n} \sum_{i=1}^n X_i, S^2 = \frac{1}{n-1} \sum_{i=1}^n (X_i - \bar{X})^2, \ T = \bar{X}^2 - \frac{1}{n} S^2$.
\begin{enumerate}
\item[(1)] 证明 $\displaystyle ET = \mu^2$;
\item[(2)] 当 $\displaystyle \mu=0, \sigma=1$ 时, 求 $\displaystyle DT$.
\end{enumerate}}
\begin{note}
  主要错的是 (2)
\end{note}
\ansat{371；【真题精选】 - 考点二：统计量的数字特征 - 5}